\section{Divisorial traces with bounded remainder}
\label{sec:resolved-trace-masses}

The invariant $\mathcal R(a;Y)$ is a supremum over all positive currents in a class. The trace of a given potential on $Y$, on the other hand, carries a specified singularity, and its non-pluripolar mass can be smaller. Suppose that on a resolution the trace becomes a divisorial logarithmic term plus a bounded remainder. The non-pluripolar product then ignores the divisor and sees only the residual current, which has bounded potentials. The class of this residual current is nef, and its top intersection number computes the mass.

We keep the assumption that $X$ belongs to class $\mathcal C$. Let $Y\subset X$ be an irreducible reduced subvariety of dimension $k>0$, let $\alpha_0$ be a smooth closed real $(1,1)$-form on $X$, and let $u\in\PSH(X,\alpha_0)$ be such that its pointwise restriction $v=u|_Y$ is not identically $-\infty$. On a K\"ahler resolution $\pi\colon Z\to Y$, with $f=j\circ\pi\colon Z\to X$ as in \cref{sec:finite-restricted-invariants}, put $\alpha_Z=f^*\alpha_0$ and $v_\pi=u\circ f$. When $Y$ is smooth, one may instead start from a smooth closed real $(1,1)$-form $\alpha_Y$ on $Y$ and $v\in\PSH(Y,\alpha_Y)$, and put $\alpha_Z=\pi^*\alpha_Y$ and $v_\pi=v\circ\pi$. In both cases $v_\pi$ is the composition of a given potential with a map; the intrinsic multiplier-ideal trace of \cref{chap:trace-data} is a different construction.

Suppose that on one such resolution
\begin{equation}\label{eq:resolved-trace-decomposition}
 v_\pi=\ell+w,
 \qquad \ell=\sum_{j=1}^N c_j\log|s_j|_{h_j}^2,
 \qquad w\text{ bounded},
\end{equation}
where $c_j>0$, the $D_j$ are prime divisors on $Z$, $s_j$ is the canonical section of $\mathcal O(D_j)$, and $h_j$ is a smooth Hermitian metric on $\mathcal O(D_j)$; the empty sum $\ell=0$ is allowed. Write $\Gamma=\sum_jc_jD_j$ and $\gamma_j=c_1(\mathcal O(D_j),h_j)$ for the curvature forms, and define the \emph{residual form} and the \emph{residual class}
\begin{equation}\label{eq:resolved-trace-residual-class}
 \eta_v=\alpha_Z-\sum_jc_j\gamma_j,
 \qquad b_v=\{\eta_v\}=\{\alpha_Z\}-\{[\Gamma]\}.
\end{equation}
For singular $Y$, all assertions below are made on its resolutions.

\begin{lemma}\label{lem:bounded-residual-mass-polynomial}\index{Residual class!mass on a resolution}
The current $R_v=\eta_v+\ddc w$ is positive and has bounded local potentials, and its class $b_v$ is nef. Moreover,
\begin{equation}\label{eq:resolved-trace-mass}
 \int_Z\langle(\alpha_Z+\ddc v_\pi)^k\rangle
       =\int_ZR_v^k
       =\int_Zb_v^k
       =\mathcal R_Z(b_v).
\end{equation}
\Needspace{7\baselineskip}
If $v_1,\ldots,v_k$ satisfy the same conditions on $Z$, possibly for different closed backgrounds, then
\begin{equation}\label{eq:resolved-trace-mixed-mass}
 \int_ZR_{v_1}\wedge\cdots\wedge R_{v_k}
       =\int_Zb_{v_1}\wedge\cdots\wedge b_{v_k}.
\end{equation}
\end{lemma}
\begin{proof}
Off the support $|\Gamma|$, we have $\ddc\ell=-\sum_jc_j\gamma_j$, so the current $\alpha_Z+\ddc v_\pi$ equals $\eta_v+\ddc w$ there and is positive. After the addition of a smooth local potential for $\eta_v$, the bounded function $w$ is plurisubharmonic off $|\Gamma|$. By \cref{lem:qpsh-removable-analytic}, it extends across the divisor as a bounded plurisubharmonic function that agrees almost everywhere with the original one; we use this extension to represent $w$. Thus $R_v\geq0$ on $Z$, and the Poincar\'e--Lelong formula gives
\[
 \alpha_Z+\ddc v_\pi=[\Gamma]+R_v.
\]
Since $w$ is bounded, $b_v$ is nef by \cref{lem:bounded-closed-potential-nef}. The non-pluripolar measure on the left of \eqref{eq:resolved-trace-mass} agrees with $R_v^k$ off the divisor, and neither measure charges the divisor. This gives the first equality.

Fix a K\"ahler form $\kappa$ on $Z$. For all sufficiently large real $t$, the smooth closed form $\eta_v+t\kappa$ is K\"ahler, and $w$ is a bounded potential for it. Therefore
\[
 \int_Z(R_v+t\kappa)^k=\int_Z(\eta_v+t\kappa)^k
\]
by \cref{thm:mass-invariance-characterization}, or by Stokes' theorem when $k=1$. Both sides are polynomials in $t$, so their constant coefficients agree; this gives $\int_ZR_v^k=\int_Zb_v^k$. The last equality in \eqref{eq:resolved-trace-mass} is \eqref{eq:nef-restricted-mass}. Applying the same polynomial argument to positive linear combinations of the $R_{v_i}$ and comparing coefficients gives \eqref{eq:resolved-trace-mixed-mass}.
\end{proof}

We denote the common value in \eqref{eq:resolved-trace-mass} by $\mathfrak r(v;Y)$ and call it the \emph{trace mass} of $v$ on $Y$. When $Y$ is smooth, it is also the non-pluripolar mass of the trace current on $Y$ itself:
\[
 \mathfrak r(v;Y)=\int_Y\langle(\alpha_Y+\ddc v)^k\rangle,
\]
where $\alpha_Y=\alpha_0|_Y$ when $v$ is a restriction. Indeed, the currents on $Z$ and on $Y$ correspond on the locus where $\pi$ is an isomorphism, and their non-pluripolar measures give no mass to its analytic complement. A change of variables therefore identifies the total masses.

\begin{theorem}\label{thm:resolved-trace-invariance}\index{Restricted volume!resolution independence}
The number $\mathfrak r(v;Y)$ does not depend on the divisor metrics or on the K\"ahler resolution carrying \eqref{eq:resolved-trace-decomposition}. If two potentials in the same closed background have the same ordinary singularity type and each has a decomposition of the form \eqref{eq:resolved-trace-decomposition}, then their residual classes agree on a common refinement, and in particular their trace masses are equal. Residual classes add when the traces and their backgrounds are added, and their mixed masses are given by \eqref{eq:resolved-trace-mixed-mass}.
\end{theorem}
\begin{proof}
Changing a divisor metric changes $\ell$ by a smooth function, $w$ by its negative, and $\eta_v$ by the corresponding $\ddc$-term. Thus $R_v$ is unchanged. The same cancellation occurs when a smooth change of the background is accompanied by the opposite change of the potential.

Let $\rho\colon W\to Z$ be a modification with $W$ K\"ahler. Pulling back the decomposition, with the total divisor $\rho^*\Gamma$ including all exceptional coefficients, gives a decomposition on $W$ with bounded remainder $w\circ\rho$, residual current $\rho^*R_v$ and nef residual class $\rho^*b_v$. Its mass is
\[
 \int_W(\rho^*b_v)^k=\int_Zb_v^k.
\]
For two decompositions on different resolutions, pass to a common K\"ahler refinement and pull both back in this way. At a generic smooth point of a prime component of the union of the two divisors, the difference of the logarithmic terms is $(a_1-a_2)\log|t|^2+O(1)$, where $t$ is a local defining coordinate of the component and $a_1,a_2$ are its two coefficients. The two potentials have bounded difference, and the resulting pair of inequalities forces $a_1=a_2$. Hence the total divisors coincide, and so do the residual classes in the common background. Finally, adding two decompositions adds both the divisors and the bounded remainders. The residual class of the sum is therefore the sum of the nef residual classes, and \eqref{eq:resolved-trace-mixed-mass} computes its mixed terms.
\end{proof}

\begin{proposition}\label{prop:resolved-trace-strict-current}
Suppose that $\mathfrak r(v;Y)>0$. Then the class $\{\alpha_Z\}$ contains a K\"ahler current $\alpha_Z+\ddc p$ whose potential satisfies $p\leq v_\pi+A$ for some constant $A$. If $b_v$ has a smooth positive representative on a resolution carrying \eqref{eq:resolved-trace-decomposition}, then on that resolution $p$ can be chosen with $p\sim v_\pi$, that is, with $p-v_\pi$ bounded. When $Y$ is smooth, the K\"ahler current and the corresponding inequalities between potentials descend to $Y$.
\end{proposition}
\begin{proof}
Positive trace mass means $\int_Zb_v^k>0$. Since $b_v$ is nef, \eqref{eq:normalized-diagonal-current} gives a current $K=\eta_v+\ddc a\geq\delta\kappa$ in $b_v$, where $\kappa$ is a K\"ahler form and $\delta>0$. The current
\[
 [\Gamma]+K=\alpha_Z+\ddc(\ell+a)
\]
is then a K\"ahler current. Since $a$ is bounded above on $Z$ and $w$ is bounded below, its potential $p=\ell+a$ satisfies $p\leq\ell+w+A=v_\pi+A$ for a constant $A$. If $b_v$ has a smooth positive representative, $a$ can be chosen smooth, and then $p-v_\pi=a-w$ is bounded.

For smooth $Y$, take $\alpha_Y=\alpha_0|_Y$ when $v$ is a restriction. By \cref{prop:psh-descent-modification}, $p=q\circ\pi$ for some $q\in\PSH(Y,\alpha_Y)$. The current $\alpha_Y+\ddc q$ is strict by \eqref{eq:strict-current-pushforward}, and since $\pi$ is surjective, each inequality between potentials descends. The hypothesis of a smooth positive representative refers to the specified resolution: the pullback of a K\"ahler class under a further blowup need not be K\"ahler.
\end{proof}

The trace mass can vanish even when the restricted class is big, because the non-pluripolar product ignores the divisorial part of the trace.

\begin{example}\label{ex:divisorial-trace-zero-mass}
Let $k\geq1$, let $Y\simeq\mathbb P^k$ be a hyperplane in $X=\mathbb P^{k+1}$, and let another hyperplane $H_X\subset X$ cut out $H=H_X\cap Y$. Let $\omega_{\mathrm{FS}}$ be the Fubini--Study form, normalized by $\int_Y\omega_{\mathrm{FS}}^k=1$, and let $s_{H_X}$ be a section of $\mathcal O(1)$ defining $H_X$, measured with the Fubini--Study metric. For $c>0$ and $d\geq0$, take $\alpha_0=(c+d)\omega_{\mathrm{FS}}$ on $X$ and $u=c\log|s_{H_X}|^2$. The restriction $v=u|_Y$ is not identically $-\infty$, and
\[
 \alpha_0|_Y+\ddc v=c[H]+d\omega_{\mathrm{FS}}|_Y,
 \qquad \mathfrak r(v;Y)=d^k,
 \qquad \mathcal R(\{\alpha_0\};Y)=(c+d)^k.
\]
Thus for $d=0$ the trace mass is zero, although the restricted invariant is positive. The divisor $c[H]$ is a nonzero part of the trace current, but the non-pluripolar top product discards it.
\end{example}

The next example shows that the inequality $p\leq v_\pi+A$ in \cref{prop:resolved-trace-strict-current} cannot in general be improved to $p\sim v_\pi$, even for a bounded trace.

\Needspace{14\baselineskip}
\begin{example}\label{ex:positive-trace-bounded-type-obstruction}
Let $\mu\colon Y\to\mathbb P^2$ be the blowup of a point, with exceptional curve $E$, let $\omega_{\mathrm{FS}}$ be the Fubini--Study form on $\mathbb P^2$ with total volume one, and put $\alpha_Y=\mu^*\omega_{\mathrm{FS}}$ and $v=0$. The decomposition \eqref{eq:resolved-trace-decomposition} holds on $Z=Y$ with no divisor, the residual class $\{\alpha_Y\}$ is nef and big, and
\[
 \mathfrak r(0;Y)=\mathcal R(\{\alpha_Y\};Y)=1.
\]
Nevertheless, there is no bounded $q\in\PSH(Y,\alpha_Y)$ for which $\alpha_Y+\ddc q\geq\delta\kappa$ with $\delta>0$ and $\kappa$ K\"ahler. Indeed, since $q$ is bounded, its restriction to $E$ is defined, and the disk criterion of \cref{lem:smooth-disk-criterion} gives
\[
 \ddc(q|_E)\geq\delta\kappa|_E,
\]
because $\alpha_Y|_E=0$. Integration over the compact curve $E$ contradicts Stokes' theorem. Thus positive trace mass gives the one-sided comparison of \cref{prop:resolved-trace-strict-current}, but a K\"ahler current in the class may need additional singularities.
\end{example}

A subvariety contained in the null locus can nevertheless have positive restricted invariant; the union in \eqref{eq:finite-restricted-null-union} is needed to recover the locus.

\begin{example}\label{ex:positive-restriction-inside-null-locus}
Let $\mu\colon X\to\mathbb P^3$ be the blowup along a line $L$, and set $a=\mu^*\{\omega_{\mathrm{FS}}\}$, where $\omega_{\mathrm{FS}}$ has total volume one; the class $a$ is nef and $\int_Xa^3=1$. Since $N_{L/\mathbb P^3}\simeq\mathcal O(1)\oplus\mathcal O(1)$, the exceptional divisor is $E=\mathbb P(N_{L/\mathbb P^3})\simeq\mathbb P^1\times\mathbb P^1$, and on $E$ the map $\mu$ is the projection to $L$. Hence $\int_Ea^2=0$, whereas a section $Y$ of $E\to L$ satisfies $\int_Ya=1$. Every irreducible subvariety not contained in $E$ has positive intersection number, because $\mu^*\omega_{\mathrm{FS}}$ is positive on its nonempty open part outside $E$. By \cref{cor:finite-restricted-null-union},
\[
 E_{\mathrm{nK}}(a)=E,
 \qquad \mathcal R(a;E)=0,
 \qquad \mathcal R(a;Y)=1.
\]
The union of the subvarieties with zero invariant is the whole null locus $E$, although the curve $Y\subset E$ has positive invariant. Positivity along one subvariety does not place it outside the null locus.
\end{example}
